\documentclass[11pt,a4paper]{ctexart}
\usepackage[margin=2.2cm]{geometry}
\usepackage{amsmath,amssymb,amsthm,mathtools,booktabs,hyperref,microtype}
\hypersetup{colorlinks=true,linkcolor=black,urlcolor=blue}
\newtheorem{theorem}{定理}
\title{\bfseries 共享素数步长的深度矩、根距离与生成函数}
\author{研究札记 · F3-MOM}
\date{2026 年 9 月 29 日}
\begin{document}
\maketitle
\begin{abstract}
固定共享根素数伸缩族中，本文在已有固定精度联合分布的基础上建立随深度增长的一致尾界。全素数乘积的单射性进一步控制极深尾，从而把指数加权分布收敛范围提升到 \(\eta<1/(q-1)\)，并得到全部固定阶矩收敛。二阶统计显式恢复根距离矩阵。最后利用“超过最大根距后仅一个坐标可继续加深”的结构，给出多元概率生成函数的有限有理表示，并证明全素数生成函数在 \(|z_i|<3^{1/(q-1)}\) 的紧多圆盘上一致收敛。状态为 \texttt{PAPER-AUDITED}，尚未 Lean 形式化。
\end{abstract}

\section{固定共享根模型}
固定 \(k,s\ge1\)，令
\[
q=3^s,\qquad b=k+s.
\]
在此前构造的固定素数伸缩族中，输出深度写成
\[
D_i=b+R_i,\qquad
R_i=v_3(d-\alpha_i(n)).
\]
根距离
\[
L_{ij}=v_3(\alpha_i-\alpha_j)
\]
有限且与 \(n\) 无关。全素数参数集合 \(\mathcal P(X)\) 满足固定形状主项
\[
|\mathcal P(X)|
=
(\kappa+o(1))
\frac{X^2}{(\log X)^a},
\]
并且每个固定深度向量的点概率趋于局部 Haar 模型 \(\mu\)。

\section{统一尾界与极深尾}
令
\[
Z=\max_iR_i.
\]
对中等深度，增长模数剩余类计数与 Selberg 上界筛给出
\[
\Pr_X(Z\ge t)
\ll
3^{-t}+X^{-3/4}
\]
当 \(3^{b+t}\le X^{1/2}\)。

对全部 \(t\)，丢弃素性后按根剩余类直接计数可得
\[
\Pr_X(Z\ge t)
\ll
(\log X)^a(3^{-t}+X^{-1}).
\]

本版新增一个极深尾改进。固定任一输出乘积 \(P_i\)。其 \(q\) 个线性素数因子严格递增；若两个参数点给出同一乘积，唯一分解恢复每个素数因子，再由前两个因子恢复 \(d,n\)。因此
\[
(n,d)\longmapsto P_i(n,d)
\]
在全素数参数集上单射。

又因为
\[
0<P_i+1\ll X^q,
\]
若
\[
R_i\ge t,
\]
则 \(P_i+1\) 是 \(3^{b+t}\) 的正倍数。于是
\[
\#\{R_i\ge t\}\ll X^q3^{-t}.
\]
除以总计数后得到
\[
\boxed{
\Pr_X(Z\ge t)
\ll
(\log X)^a
\min\{3^{-t}+X^{-1},\,X^{q-2}3^{-t}\}.
}
\]

\section{指数加权分布}
固定
\[
0<\theta<\frac1{q-1}.
\]
把尾求和分成三段：
\[
t\le \tfrac12\log_3X,
\qquad
\tfrac12\log_3X<t\le(q-1)\log_3X,
\qquad
t>(q-1)\log_3X.
\]
前两段使用中等/全深度尾界，第三段使用新的单射尾界。得到
\[
\sup_X\mathbb E_X3^{\theta Z}<\infty.
\]

于是对任意
\[
0<\eta<1/(q-1)
\]
选 \(\eta<\theta<1/(q-1)\)，可在盒外用
\[
3^{\eta Z}
\le
3^{-(\theta-\eta)K}3^{\theta Z}
\]
做统一可积控制。有限盒内使用固定点概率收敛，得到
\[
\boxed{
\sum_{\mathbf r}
3^{\eta\max_i r_i}
|\mu_X(\mathbf r)-\mu(\mathbf r)|
\to0.
}
\]
所以全部固定阶多项式矩收敛。

\section{协方差恢复根距离}
局部边缘分布满足
\[
\Pr(R_i=0)=\frac12,\qquad
\Pr(R_i=t)=3^{-t}\quad(t\ge1),
\]
故
\[
\mathbb ER_i=\frac34,\qquad
\operatorname{Var}(R_i)=\frac{15}{16}.
\]

若 \(L=L_{ij}\)，当两坐标不等时较小者恰为 \(L\)，并且
\[
\Pr(R_i=L+t,R_j=L)
=
\Pr(R_i=L,R_j=L+t)
=
3^{-(L+t)}.
\]
所以
\[
\boxed{
\mathbb E(R_i-R_j)^2
=
3^{1-L_{ij}}.
}
\]
进而
\[
\boxed{
\operatorname{Cov}(D_i,D_j)
=
\frac{15}{16}
-
\frac32\,3^{-L_{ij}},
}
\]
\[
\boxed{
\operatorname{Corr}(D_i,D_j)
=
1-\frac85\,3^{-L_{ij}}.
}
\]
指数矩收敛把这些公式转移到全素数族。

\section{有限根树与多元生成函数}
令
\[
K=\max_{i\ne j}L_{ij}.
\]
若某个 \(R_i=u>K\)，则对所有 \(j\ne i\)
\[
R_j=L_{ij}.
\]
因此深层支撑只由 \(m\) 条互不相交的射线组成。

定义
\[
\Phi(\mathbf z)=\mathbb E_\mu\prod_i z_i^{R_i}.
\]
有限盒内的部分记为 \(A_K(\mathbf z)\)。在 \(|z_i|<3\) 内，
\[
\boxed{
\Phi(\mathbf z)
=
A_K(\mathbf z)
+
\sum_i
\left(\prod_{j\ne i}z_j^{L_{ij}}\right)
\frac{(z_i/3)^{K+1}}{1-z_i/3}.
}
\]
所以 \(\Phi\) 是有理函数，分母整除
\[
\prod_i(1-z_i/3).
\]

令
\[
S_0=\sum_iR_i,\qquad
s_i=\sum_{j\ne i}L_{ij}.
\]
当
\[
r>mK
\]
时，只有深层射线能贡献，因此
\[
\boxed{
\mu(S_0=r)
=
3^{-r}\sum_i3^{s_i}.
}
\]

\section{全素数生成函数收敛}
定义
\[
\Phi_X(\mathbf z)=\mathbb E_X\prod_i z_i^{R_i}.
\]
由于深层支撑满足
\[
\sum_iR_i\le Z+(m-1)K,
\]
任意紧多圆盘
\[
|z_i|\le R<3^{1/(q-1)}
\]
上有
\[
|\mathbf z^{\mathbf r}|
\le
R^{(m-1)K}R^Z.
\]
令
\[
\eta=\log_3R<1/(q-1),
\]
再用指数加权总变差，即得
\[
\boxed{
\Phi_X\to\Phi
}
\]
在这些紧集上一致收敛。固定阶复偏导数则由稍大多圆盘上的 Cauchy 积分公式得到一致收敛。

\section{边界}
局部函数在 \(|z_i|<3\) 内绝对收敛；目前全素数平均只证明到 \(|z_i|<3^{1/(q-1)}\)。本文不声称后一范围最优，也不提供增长深度事件的相对素数主项。

\begin{thebibliography}{9}
\bibitem{GT} B. Green and T. Tao, \emph{Linear equations in primes}, Ann. of Math. 171 (2010).
\bibitem{Kedlaya} K. Kedlaya, \emph{Notes on analytic number theory}, Selberg sieve chapters.
\bibitem{Zuniga} W. A. Zúñiga-Galindo, work on Igusa local zeta functions of univariate polynomials.
\end{thebibliography}
\end{document}

% final-proof-pdf-build
